% Part: first-order-logic
% Chapter: axiomatic-deduction
% Section: deduction-theorem

% verification of properties of provability needed for maximally
% consistent sets in the completeness chapter.

\documentclass[../../../include/open-logic-section]{subfiles}

\begin{document}

\iftag{FOL}
      {\olfileid[fr]{fol}{axd}{ded}}
      {\olfileid[fr]{pl}{axd}{ded}}

\olsection{Le théorème de la déduction}

Comme nous l'avons vu, donner des !!{derivation}s dans un système
axiomatique est laborieux, et celles-ci peuvent être difficiles à
trouver. Plutôt que d'écrire effectivement de longues listes
de !!{formula}s, il est généralement plus facile d'établir l'existence
de ces !!{derivation}s au moyen de quelques résultats simples. Nous en
avons déjà démontré trois. La \olref[ptn]{prop:reflexivity} permet
toujours d'affirmer $\Gamma \Proves !A$ lorsque $!A \in \Gamma$.
La \olref[ptn]{prop:monotonicity} affirme que, si
$\Gamma \Proves !A$, alors $\Gamma \cup \{!B\} \Proves !A$.
Enfin, la \olref[ptn]{prop:transitivity} entraîne que, si
$\Gamma \Proves !A$ et $!A \Proves !B$, alors
$\Gamma \Proves !B$. Voici un autre résultat simple, une version
« méta » du modus ponens :

\begin{prop}
\ollabel{prop:mp}
Si $\Gamma \Proves !A$ et $\Gamma \Proves !A \lif !B$, alors
$\Gamma \Proves !B$.
\end{prop}

\begin{proof}
Nous avons $\{!A, !A \lif !B\} \Proves !B$ :
\begin{derivation}
  1. & $!A$ & \Hyp\\
  2. & $!A \lif !B$ & \Hyp\\
  3. & $!B$ & 1, 2, \MP
\end{derivation}
La \olref[ptn]{prop:transitivity} donne alors
$\Gamma \Proves !B$.
\end{proof}

Le résultat le plus important que nous utiliserons dans ce contexte
est le théorème de la déduction.

\begin{thm}[Théorème de la déduction]
\ollabel{thm:deduction-thm}
$\Gamma \cup \{!A\} \Proves !B$ si et seulement si
$\Gamma \Proves !A \lif !B$.
\end{thm}

\begin{proof}
Le sens « si » est immédiat. Si $\Gamma \Proves !A \lif !B$, alors
$\Gamma \cup \{!A\} \Proves !A \lif !B$ par la
\olref[ptn]{prop:monotonicity}. En outre,
$\Gamma \cup \{!A\} \Proves !A$ par la
\olref[ptn]{prop:reflexivity}. La \olref{prop:mp} donne donc
$\Gamma \cup \{!A\} \Proves !B$.

Pour le sens « seulement si », procédons par récurrence sur la longueur
de la !!{derivation} de~$!B$ à partir de $\Gamma \cup \{!A\}$.

Dans le cas de base, démontrons l'assertion pour toute !!{derivation}
de longueur~$1$. Une telle !!{derivation} de~$!B$ à partir de
$\Gamma \cup \{!A\}$ se réduit à $!B$; si elle est correcte, alors
$!B \in \Gamma \cup \{!A\}$ ou bien $!B$ est un axiome. Si
$!B \in \Gamma$ ou si $!B$ est un axiome, alors
$\Gamma \Proves !B$. Nous avons aussi
$\Gamma \Proves !B \lif (!A \lif !B)$ d'après
l'\olref[prp]{ax:lif1}; la \olref{prop:mp} donne donc
$\Gamma \Proves !A \lif !B$. Si $!B \in \{!A\}$, alors
$\Gamma \Proves !A \lif !B$, car $!A \lif !B$ est dans ce cas le
même !!{sentence} que $!A \lif !A$, dont nous avons donné une
!!{derivation} dans l'\olref[pro]{ex:identity}.

Pour l'étape de récurrence, supposons qu'une !!{derivation} de~$!B$ à
partir de $\Gamma \cup \{!A\}$ se termine par une étape~$!B$ justifiée
par modus ponens. Si ce n'est pas le cas, alors $!B \in \Gamma$,
$!B \ident !A$, ou $!B$ est un axiome, et le raisonnement du cas de
base s'applique. Dans le cas du modus ponens, des étapes antérieures de
la !!{derivation} sont $!C \lif !B$ et $!C$, pour une certaine
!!{formula}~$!C$. Ainsi,
$\Gamma \cup \{!A\} \Proves !C \lif !B$ et
$\Gamma \cup \{!A\} \Proves !C$. Les !!{derivation}s correspondantes
sont plus courtes; l'hypothèse de récurrence leur est donc applicable.
Nous obtenons :
\begin{align*}
  & \Gamma \Proves !A \lif (!C \lif !B); \\
  & \Gamma \Proves !A \lif !C.
\end{align*}
Mais nous avons également
\[
\Gamma \Proves (!A \lif (!C \lif !B)) \lif
((!A \lif !C) \lif (!A \lif !B))
\]
d'après l'\olref[prp]{ax:lif2}. Deux applications de la
\olref{prop:mp} donnent donc $\Gamma \Proves !A \lif !B$, comme il
fallait le démontrer.
\end{proof}

Remarquons que les deux schémas d'axiomes correspondants,
l'\olref[prp]{ax:lif1} et l'\olref[prp]{ax:lif2}, ont été choisis
précisément de façon que le théorème de la déduction soit valable.

Les faits suivants sur la !!{derivability} sont utiles; nous les
laissons en exercices.

\begin{prop}
\ollabel{prop:derivfacts}
\begin{enumerate}
\item $\Proves (!A \lif !B) \lif ((!B \lif !C)
  \lif (!A \lif !C))$; \ollabel{derivfacts:a}
\item Si $\Gamma \cup \{\lnot !A\} \Proves \lnot !B$, alors
  $\Gamma \cup \{!B\} \Proves !A$ (contraposition);
  \ollabel{derivfacts:b}
\item $\{!A, \lnot !A\} \Proves !B$ (\emph{ex falso quodlibet},
  explosion); \ollabel{derivfacts:c}
\item $\{\lnot\lnot !A\} \Proves !A$ (élimination de la double
  négation); \ollabel{derivfacts:d}
\item Si $\Gamma \Proves \lnot\lnot !A$, alors
  $\Gamma \Proves !A$; \ollabel{derivfacts:e}
\end{enumerate}
\end{prop}

\tagprob{FOL}
\begin{prob}
Démontrer la \olref[fol][axd][ded]{prop:derivfacts}.
\end{prob}
\tagendprob

\tagprob{notFOL}
\begin{prob}
Démontrer la \olref[pl][axd][ded]{prop:derivfacts}.
\end{prob}
\tagendprob

\begin{editorial}
Dans le premier item de la proposition précédente, la source omet la
parenthèse fermante du conséquent. Elle a été rétablie.
\iftag{FOL}{En présence des quantificateurs, l'étape supplémentaire
qui traite les règles~\QR{} est donnée dans la section suivante.}{}
\end{editorial}

\end{document}
